% Build from the repository root:
% latexmk -pdf -interaction=nonstopmode -halt-on-error \
%   -outdir=tmp/pdfs/uop-algebra paper/uop-algebra.tex
% The Haskell listings form one module when preceded by:
% module UOpExamples where
% import Control.Applicative ((<|>))
% import Data.List (sortOn, transpose)
% import Data.Maybe (fromMaybe)
% import Numeric.Natural (Natural)
\documentclass[11pt,a4paper]{article}
\usepackage[hmargin=23mm,vmargin=22mm,headheight=14pt,headsep=7mm,footskip=11mm]{geometry}
\usepackage[T1]{fontenc}
\usepackage{amsmath,amsthm}
\usepackage{newtxtext,newtxmath}
\usepackage{microtype}
\usepackage{booktabs,tabularx,array}
\usepackage{enumitem,needspace,etoolbox}
\usepackage{xcolor,listings,fancyhdr}
\usepackage{hyperref}
\definecolor{ink}{HTML}{173A54}
\definecolor{codebg}{HTML}{F3F5F7}
\hypersetup{colorlinks=true,linkcolor=ink,citecolor=ink,urlcolor=ink,
  pdftitle={Algebra and Rewriting for UOp},
  pdfsubject={Matrix identities and simple Haskell rewrites for a list-based unitary IR}}
\pagestyle{fancy}
\fancyhf{}
\fancyhead[L]{\small Algebra and rewriting for UOp}
\fancyhead[R]{\small QPP / HQP}
\fancyfoot[C]{\thepage}
\renewcommand{\headrulewidth}{0.3pt}
\setlength{\parindent}{0pt}
\setlength{\parskip}{4pt plus 1pt}
\setlength{\emergencystretch}{2em}
\setlist{nosep,leftmargin=1.5em}
\lstset{language=Haskell,basicstyle=\small\ttfamily,
  backgroundcolor=\color{codebg},frame=single,rulecolor=\color{black!15},
  framerule=0.3pt,framesep=5pt,columns=fullflexible,
  keepspaces=true,showstringspaces=false,breaklines=true,
  aboveskip=7pt,belowskip=7pt}
\BeforeBeginEnvironment{lstlisting}{\par\noindent\begin{minipage}{\linewidth}}
\AfterEndEnvironment{lstlisting}{\end{minipage}\par}
\AddToHook{cmd/section/before}{\Needspace{8\baselineskip}}
\newcommand{\C}{\mathbb C}
\newcommand{\N}{\mathbb N}
\newcommand{\Q}{\mathbb Q}
\newcommand{\sem}[1]{\llbracket #1\rrbracket}
\newcommand{\code}[1]{\texttt{#1}}
\newcommand{\Id}{I}
\newcommand{\Rot}[2]{R_{#1}(#2)}
\newcommand{\Ph}[1]{\phi(#1)}
\newcommand{\ds}{\boldsymbol d}
\DeclareMathOperator{\diag}{diag}
\newtheorem{proposition}{Proposition}
\newcolumntype{L}{>{\raggedright\arraybackslash}X}

\begin{document}
\begin{center}
{\LARGE\bfseries Algebra and Rewriting for UOp}\\[8pt]
{\large Equational laws, rewrite patterns, and transformation strategies}\\[7pt]
{\small QPP / HQP \quad$\cdot$\quad 17 October 2026}
\end{center}

\begin{abstract}
\noindent Flattened operand lists expose algebraic rewrites as maps,
guarded zips, identity-padded transposes, adjacent-pair scans, and list
splits. This note derives their exact, phase-preserving identities and
gives a small Haskell implementation. Permutations account for basis
order; strategies choose rule direction, traversal, and growth limits.
A QFT example shows how lifting compositions produces parallel layers
while increasing the number of identity nodes. All transformations
operate symbolically.
\end{abstract}

\section{Syntax, semantics, and QOp lowering}

A well-formed UOp denotes a unitary endomorphism of $\C^d$, for any $d\in\N$,
including $d=0$. Dimensions count basis vectors, not qubits. The space
$\C^0=\{0\}$ has a unique endomorphism, the unitary $I_0$.
The current core is:

\begin{lstlisting}
type Dim = Natural
type Permutation = [Int]

data UOp
  = R UOp Rational
  | Tensor [UOp]
  | DirectSum [UOp]
  | Compose Dim [UOp]
  | PermuteTensor [Dim] Permutation
  | PermuteSum [Dim] Permutation
  deriving (Show, Eq)
\end{lstlisting}

\code{Natural} is unbounded. A permutation must contain each index
$0,\ldots,k-1$ exactly once, where $k$ is the length of its dimension
list. Its convention is \emph{output position $j$ selects input position
$p_j$}. Dimensions are needed because the same index permutation can
denote different matrices for different block sizes.

Write $A=\sem a$ and $d(a)=\dim A$. We write $a\equiv b$ when the
terms have the same dimension and the same matrix denotation, including
global phase. Composition is in mathematical order:
\code{Compose d [a,b]} denotes $AB$, applying $b$ first.
Every operand of \code{Compose d xs} must have dimension $d$.
The dimension annotation also determines the empty composition.

The rotation uses the same sign and half-angle convention as QOp:
\[
  \Rot A t:=e^{-i\pi tA/2},\qquad
  \sem{\code{R }a\ t}=\Rot A t,\qquad t\in\Q.
\]
For $t\ne0$, its already-unitary axis must be proved Hermitian.
At $t=0$, the axis must still be a valid UOp but need not be Hermitian.
Thus a nonzero-angle axis satisfies $A^2=I$.
Lowering QOp's \code{R a t} keeps the angle $t$ unchanged and recursively
lowers its axis. These are semantic well-formedness conditions; the raw
Haskell datatype does not enforce them.

\clearpage
\subsection{Deriving QOp constructs from UOp}
\label{sec:qop-lowering}

The following are semantic identities, including global phases. Write
$\widehat a$ for the recursively translated QOp $a$, and
$d=\dim\widehat a$. All source terms are assumed well formed.
In the right column, $I_d$ denotes the identity on $\C^d$,
$\Ph q$ abbreviates the scalar phase, and $x,z$ abbreviate the UOp
expressions defined in their rows. Infix $\otimes$ and $\circ$ stand
for \code{Tensor [a,b]} and \code{Compose d [a,b]}, respectively;
$R_a(t)$ stands for \code{R a t}.

\begin{center}
\begingroup
\renewcommand{\arraystretch}{1.18}
\begin{tabularx}{\linewidth}{@{}lL@{}}
\toprule
Identity or QOp construct & Equivalent UOp expression\\
\midrule
$I_0$ & \code{Compose 0 []} $\equiv$ \code{DirectSum []}\\
$I_1$ & \code{Compose 1 []} $\equiv$ \code{Tensor []}\\
$I_2$ & \code{Compose 2 []} $\equiv$ $\code{DirectSum }[I_1,I_1]$\\
$I_d$ & \code{Compose d []}\\
\midrule
\code{Id n} & $\code{Compose }(2^n)\,[]\equiv I_{2^n}$\\
\code{Phase q} & $\Ph q:=\code{R }I_1\,(-2q)$\\
\code{X} & $x:=\code{PermuteSum }[1,1]\,[1,0]$\\
\code{Z} & $z:=\code{DirectSum }[I_1,\Ph1]$\\
\code{Y} & $\Ph{1/2}\otimes(x\circ z)$\\
\code{H} & $\Ph{1/2}\otimes
  \bigl(R_z(1/2)\circ R_x(1/2)\circ R_z(1/2)\bigr)$\\
\code{SX} & $\Ph{1/4}\otimes R_x(1/2)$\\
\code{R a t} & $\code{R }\widehat a\ t$\\
\code{C a} & $\code{DirectSum }[I_d,\widehat a]$\\
\code{Permute p} & $\code{PermuteTensor }(\code{replicate }|p|\ 2)\ p$\\
\code{Tensor a b} & $\code{Tensor }[\widehat a,\widehat b]$\\
\code{DirectSum a b} & $\code{DirectSum }[\widehat a,\widehat b]$\\
\code{Compose a b} & $\code{Compose }d\,[\widehat a,\widehat b]$\\
\code{Adjoint a} & $\widehat a^{\dagger}$ (structural adjoint; Section~\ref{sec:adjoints})\\
\bottomrule
\end{tabularx}
\endgroup
\end{center}

\paragraph{Dimension versus qubit count.}
QOp's \code{One = Id 0} translates to $I_1$, and
\code{I = Id 1} translates to $I_2$. The zero-dimensional $I_0$ is a
valid UOp but has no QOp counterpart. The preferred representation of
\code{Id n}, also used by the implementation, is
\code{Compose }$(2^n)$\code{ []}: one constructor with an empty operand
list. Its syntax-tree size is constant; the dimension's binary encoding
uses $n+1$ bits. The equivalent \code{Tensor }$(\code{replicate }n\ I_2)$
uses $n$ factors and is useful when the qubit decomposition should be
explicit. Neither representation enumerates $2^n$ scalar blocks.

\paragraph{Phase and gate conventions.}
The phase parameter must be $-2q$, since
$e^{-i\pi(-2q)I_1/2}=e^{i\pi q}I_1$.
The gate identities use $Y=iXZ$ and
$R_Z(1/2)R_X(1/2)R_Z(1/2)=-iH$.
The chosen square root is
\[
 \Ph{1/4}\otimes R_X(1/2)
 =\frac{1+i}{2}I_2+\frac{1-i}{2}X,
 \qquad (\mathrm{SX})^2=X.
\]
Thus the factors $\Ph{1/2}$ and $\Ph{1/4}$ are necessary for exact
equality, particularly when a gate is placed inside a controlled block.
The abbreviations and structural adjoint add no constructors to UOp.

\section{Rewrite contracts and structural normalization}

\label{sec:contracts}

The examples operate on well-formed terms. Establish validity at the
entry point by validation or a construction known to preserve it;
internal rewrites then preserve it by their stated premises and
equality proofs. This can be maintained as an internal invariant.
Explicit proof objects are an optional implementation technique.
An incomplete Hermiticity checker may reject a rewritten axis even when
it is equal to the original; a trusted rewrite preserves the property
without requiring the checker to rediscover it.

The listings form a small reference implementation of these patterns.
They assume valid input and are separate from the repository's public
API; Section~\ref{sec:implementation} records what is implemented there.

A local rewrite $a\longrightarrow b$ must establish $a\equiv b$ and
well-formedness of the replacement. Its contract has four parts:
\begin{enumerate}
\item \textbf{Match:} which constructor or list fragment is recognized?
\item \textbf{Premises:} which dimensions, permutations, Hermiticity
facts, or syntactic equalities make the equation applicable?
\item \textbf{Replacement:} what term represents the other side?
\item \textbf{Direction and cost:} why use this orientation in this pass,
and can it duplicate subterms or introduce additional nodes?
\end{enumerate}

For example, fusion matches the same rotation axis, and permutation
movement matches the full factor or block decomposition. Constructor
names alone do not establish these premises.

\begin{center}
\begin{tabularx}{\linewidth}{@{}lLL@{}}
\toprule
Pattern & Algebraic use & Haskell operation\\
\midrule
All children & Adjoint, block rotation, traversal & \code{map}\\
Aligned lists & Compose matching tensor factors or sum blocks & Guard dimensions, then \code{zipWith}\\
Unequal product lengths & Lift composition out of parallel operands & Pad with identities, then \code{transpose}\\
Adjacent factors & Inverse cancellation, rotation/permutation fusion & One shared list scan\\
Selected operand & Distribute a tensor over one direct sum & \code{break}, retain prefix and suffix\\
Whole tree & Select locations and repeat a directed rule set & Bottom-up traversal with an explicit stopping policy\\
\bottomrule
\end{tabularx}
\end{center}

\paragraph{Equality is stronger than equality up to phase.}
For a phase $\lambda$, $\lambda A$ may give the same isolated quantum
state up to phase, but $I\oplus\lambda A$ generally differs observably
from $I\oplus A$. Rewrites must therefore preserve exact operator
equality in every context. The scalar-phase rules below do so.

\emph{Well-formedness} establishes validity; \emph{normalization} fixes
representational invariants; \emph{optimization} improves a stated cost.
For example, lifting compositions reduces nesting but can increase size.

\subsection{Associativity, units, and dimension zero}

Use the usual lexicographic basis for tensor products, with the first
factor varying slowest, and concatenated bases for direct sums. Then
associativity holds as literal matrix equality:
\[
 (A\otimes B)\otimes C=A\otimes(B\otimes C),\qquad
 (A\oplus B)\oplus C=A\oplus(B\oplus C),\qquad
 (AB)C=A(BC).
\]
Hence nested nodes of the same kind can be flattened without a
permutation. For example, \code{Tensor [a,Tensor [b,c],d]} becomes
\code{Tensor [a,b,c,d]}. Operand order is preserved.

\begin{center}
\begin{tabularx}{\linewidth}{@{}lLL@{}}
\toprule
List operation & Dimension & Empty list / unit\\
\midrule
$\bigotimes_i A_i$ & $\prod_i d_i$ & $I_1$\\
$\bigoplus_i A_i$ & $\sum_i d_i$ & $I_0$\\
$A_1\cdots A_k$ & Common dimension $d$ & $I_d$\\
\bottomrule
\end{tabularx}
\end{center}
Every singleton list denotes its sole operand. The unit laws are
\[
 I_1\otimes A=A=A\otimes I_1,\qquad
 I_0\oplus A=A=A\oplus I_0,\qquad
 I_dA=A=AI_d\quad(\dim A=d).
\]
Only these units may be deleted. In particular, $I_2\otimes A$ and
$I_2\oplus A$ generally have different dimensions from $A$.
Products and sums of identities can instead be combined:
\[
 \bigotimes_i I_{d_i}=I_{\prod_i d_i},\qquad
 \bigoplus_i I_{d_i}=I_{\sum_i d_i},\qquad
 I_0\otimes A=I_0=A\otimes I_0.
\]
The last law follows because a tensor product with a zero-dimensional
space has dimension zero. Every valid dimension-zero term can therefore
be represented by the same canonical identity.

\begin{lstlisting}
identity :: Dim -> UOp
identity d = Compose d []

dim :: UOp -> Dim
dim (R a _)             = dim a
dim (Tensor xs)          = product (map dim xs)
dim (DirectSum xs)       = sum (map dim xs)
dim (Compose d _)        = d
dim (PermuteTensor ds _) = product ds
dim (PermuteSum ds _)    = sum ds

isIdentity :: UOp -> Bool
isIdentity (Compose _ []) = True
isIdentity _             = False
\end{lstlisting}

\code{isIdentity} recognizes the chosen syntax for identities; it does
not decide whether an arbitrary term denotes an identity. The same
distinction applies to the derived \code{Eq} instance: syntactic equality
is a sufficient, incomplete test for semantic equality.

\Needspace{34\baselineskip}
\subsection{Structural builders}

These internal builders take valid operands with structurally normalized
children. In particular, \code{compose d xs} requires $\code{dim x}=d$
for every operand. Each builder flattens its own operation and removes
its unit; \code{finish} handles dimension zero, identities, and singletons.

\begin{lstlisting}
tensor :: [UOp] -> UOp
tensor xs = finish Tensor (product (map dim ys)) ys
  where
    ys = filter (/= identity 1) (concatMap unpack xs)
    unpack (Tensor zs) = zs
    unpack x           = [x]

directSum :: [UOp] -> UOp
directSum xs = finish DirectSum (sum (map dim ys)) ys
  where
    ys = filter (/= identity 0) (concatMap unpack xs)
    unpack (DirectSum zs) = zs
    unpack x              = [x]

compose :: Dim -> [UOp] -> UOp
compose d xs = finish (Compose d) d (concatMap unpack xs)
  where
    unpack (Compose _ zs) = zs  -- Empty compositions disappear.
    unpack x              = [x]

finish :: ([UOp] -> UOp) -> Dim -> [UOp] -> UOp
finish con d xs
  | d == 0 || all isIdentity xs = identity d
  | [x] <- xs                  = x
  | otherwise                  = con xs
\end{lstlisting}

Since \code{all isIdentity []} is true, \code{finish} handles all
three empty lists with their respective dimensions.

For example, assuming $a$ is normalized,
\[
\begin{aligned}
 \code{tensor [identity 1,a]}&=a,\\
 \code{directSum [identity 2,identity 3]}&=\code{identity 5},\\
 \code{tensor [identity 2,identity 3]}&=\code{identity 6},\\
 \code{compose (dim a) [identity (dim a),a]}&=a.
\end{aligned}
\]
These are structural invariants, not a complete normal form for matrix
equality. For example, adjacent identity factors inside a larger mixed
tensor list are not combined by this minimal implementation.

\paragraph{Validate before erasing.}
Check raw input recursively before normalization: erasing a zero-dimensional
tensor could otherwise hide an invalid child, and flattening a composition
could discard an incorrect dimension annotation. Internal builders may
rely on the validity contract established at the entry point.

\section{Permutations and adjoints}

Let $\ds=(d_0,\ldots,d_{k-1})$ and
$\ds_p=(d_{p_0},\ldots,d_{p_{k-1}})$.
The tensor permutation $T(\ds,p)$ sends
\[
 v_0\otimes\cdots\otimes v_{k-1}
 \longmapsto v_{p_0}\otimes\cdots\otimes v_{p_{k-1}}.
\]
The sum permutation $S(\ds,p)$ sends a concatenation of vectors to
the concatenation in that same reordered block order. The matrices are
square, of dimensions $\prod_i d_i$ and $\sum_i d_i$, respectively.
Their input and output \emph{decompositions} can nevertheless differ.

For either kind of permutation, writing $P$ for $T$ or $S$,
\begin{align}
 P(\ds,\mathrm{id})&=I, \label{eq:perm-id}\\
 P(\ds,p)^\dagger&=P(\ds_p,p^{-1}), \label{eq:perm-adjoint}\\
 P(\ds_p,q)P(\ds,p)&=P(\ds,r),
       \qquad r_j=p_{q_j}. \label{eq:perm-compose}
\end{align}
Follow a factor or block through the reorderings to obtain these laws.
The adjoint needs both the inverse index permutation and the reordered
dimension list.

For example, $S([1,2],[1,0])$ maps $(x,y,z)$ to $(y,z,x)$.
It is a three-cycle, although $[1,0]$ is an involution. Its adjoint is
$S([2,1],[1,0])$, not another copy of itself.
Therefore an involutive \emph{block-index permutation} alone does not
establish Hermiticity when block sizes differ.

For $\dim A_i=d_i$, the naturality laws are
\begin{align}
 T(\ds,p)\left(\bigotimes_i A_i\right)T(\ds,p)^\dagger
    &=\bigotimes_j A_{p_j}, \label{eq:tensor-naturality}\\
 S(\ds,p)\left(\bigoplus_i A_i\right)S(\ds,p)^\dagger
    &=\bigoplus_j A_{p_j}. \label{eq:sum-naturality}
\end{align}
For tensors, apply both sides to a pure tensor; these span the space.
For sums, apply both sides to a vector supported in one block.
These equations justify reordering \emph{with} conjugating permutations.
Neither tensor nor direct sum is generally commutative as a matrix
operation, so sorting their operand lists is not a valid normalization.

\paragraph{Associativity does not make permutation metadata optional.}
Flattening a tensor or sum preserves its matrix. But a permutation that
refers to the old grouping must retain those grouped dimensions, or
be explicitly translated to the new grouping. Likewise, removing a
dimension-one tensor factor or a dimension-zero sum block requires
updating any permutation that refers to its position. The list index
alone is not a stable identifier for an operand.

\subsection{Helpers and structural adjoints}

\label{sec:adjoints}

For readability, use lists for the small permutation metadata. The
following helper interprets our output-selects-input convention directly:

\begin{lstlisting}
reorder :: Permutation -> [a] -> [a]
reorder p xs = map (xs !!) p

inverse :: Permutation -> Permutation
inverse p = map snd (sortOn fst (zip p [0..]))

permutation :: ([Dim] -> Permutation -> UOp)
            -> ([Dim] -> Dim) -> [Dim] -> Permutation -> UOp
permutation con size ds p
  | d <= 1 || p == [0 .. length ds - 1] = identity d
  | otherwise                         = con ds p
  where d = size ds

permuteTensor, permuteSum :: [Dim] -> Permutation -> UOp
permuteTensor = permutation PermuteTensor product
permuteSum    = permutation PermuteSum sum
\end{lstlisting}

These helpers require valid permutations. In dimensions zero and one,
every permutation matrix is the identity. If metadata becomes large,
replace list indexing in \code{reorder} with an array or vector; no
basis permutation of size $\prod_i d_i$ need be materialized.

Adjoints are derived operations, so a separate constructor is unnecessary:
\[
 (AB)^\dagger=B^\dagger A^\dagger,\quad
 (A\otimes B)^\dagger=A^\dagger\otimes B^\dagger,\quad
 (A\oplus B)^\dagger=A^\dagger\oplus B^\dagger.
\]
Only composition reverses its list. For accepted rotations,
$\Rot A t^\dagger=\Rot A{-t}$; the axis remains unchanged.

\begin{lstlisting}
adjoint :: UOp -> UOp
adjoint op = case op of
  R a t            -> rotation a (-t)
  Tensor xs        -> tensor (map adjoint xs)
  DirectSum xs     -> directSum (map adjoint xs)
  Compose d xs     -> compose d (map adjoint (reverse xs))
  PermuteTensor ds p ->
    permuteTensor (reorder p ds) (inverse p)
  PermuteSum ds p ->
    permuteSum (reorder p ds) (inverse p)
\end{lstlisting}

The rotation helper is given next. All well-formed terms are unitary, so
$A^\dagger=A^{-1}$ and $AA^\dagger=A^\dagger A=I_{\dim A}$.
These laws permit adjacent cancellation in a composition list; they do
not permit cancelling separated operands across an arbitrary intervening
operator. Applying \code{adjoint} twice preserves the denotation, although
normalization may change the syntax.

\section{Phases and rotations}

\label{sec:phases}

Scalar phase is an abbreviation:
\[
 \Ph t=[e^{i\pi t}]=\Rot{I_1}{-2t},
 \qquad \operatorname{scale}(t,A)=\Ph t\otimes A=e^{i\pi t}A.
\]
Thus \code{Phase} is optional in the core. Its derived implementation
costs one rotation node. Scalar rules include
\[
 \Ph0=I_1,\quad \Ph{s+2k}=\Ph s,\quad
 \Ph s\Ph t=\Ph s\otimes\Ph t=\Ph{s+t},\quad
 \Ph t^\dagger=\Ph{-t}.
\]
Scalar factors commute with every operator after lifting to its
dimension. Under composition and tensor product their phases add:
\[
 \operatorname{scale}(s,A)\operatorname{scale}(t,B)
   =\operatorname{scale}(s+t,AB),\qquad
 \operatorname{scale}(s,A)\otimes\operatorname{scale}(t,B)
   =\operatorname{scale}(s+t,A\otimes B).
\]
A \emph{common} phase may move across a direct sum:
$\bigoplus_i\operatorname{scale}(t,A_i)
=\operatorname{scale}(t,\bigoplus_i A_i)$.
Different block phases remain relative phases; they cannot simply be added.
In a flat composition a scalar phase is attached to \emph{one} factor:
\[
 \operatorname{scale}(t,A_1\cdots A_m)
 =\operatorname{scale}(t,A_1)A_2\cdots A_m\qquad(m>0).
\]
Attaching it to every factor would multiply the phase by $m$. Conversely,
scaling a direct sum must scale \emph{every} block. These two cases must
remain distinct in the implementation.

If $A$ is unitary and Hermitian, $A^2=I$. Splitting the exponential's
power series into its even and odd terms gives
\begin{equation}
 \Rot A t=\cos(\pi t/2)I-i\sin(\pi t/2)A. \label{eq:involution-exp}
\end{equation}
In particular, for $k\in\mathbb Z$,
\[
 \Rot A0=I,\quad \Rot A{t+4k}=\Rot At,\quad
 \Rot A{2k}=(-1)^kI,\quad \Rot A{2k+1}=-i(-1)^kA.
\]
Also $R_{-A}(t)=R_A(-t)$ and
$R_{I_d}(t)=\operatorname{scale}(-t/2,I_d)$.
The following minimal helper implements only zero, periodicity, and
dimension zero. Special integer angles can be a separate rule.

\begin{lstlisting}
modFour :: Rational -> Rational
modFour t = t - 4 * fromInteger (ceiling ((t - 2) / 4))

rotation :: UOp -> Rational -> UOp
rotation a t
  | dim a == 0 || r == 0 = identity (dim a)
  | otherwise           = R a r
  where r = modFour t

phase :: Rational -> UOp
phase t = rotation (identity 1) (-2*t)

scale :: Rational -> UOp -> UOp
scale t a = tensor [phase t, a]
\end{lstlisting}

\textbf{The precondition concerns the original angle.}
For $t\ne0$, \code{rotation} requires a proved Hermitian axis before
reducing $t$ modulo four. For example, a non-Hermitian axis with angle
$4$ must not be accepted merely because $4\bmod4=0$.
At the original angle zero, a valid non-Hermitian axis is allowed.

\subsection{Rotation rewrites}

For compatible dimensions and accepted rotations, the useful laws are
\begin{align}
 \Rot As\Rot At&=\Rot A{s+t}, \label{eq:exp-fuse}\\
 U\Rot At U^\dagger&=\Rot{UAU^\dagger}t, \label{eq:exp-conjugate}\\
 \Rot{\bigoplus_i A_i}t&=\bigoplus_i\Rot{A_i}t,
       \label{eq:exp-sum}\\
 \Rot{I_a\otimes A\otimes I_c}t
     &=I_a\otimes\Rot At\otimes I_c. \label{eq:exp-id-tensor}
\end{align}
Equation~\eqref{eq:exp-fuse} follows because scalar multiples of the same
matrix commute. For conjugation, $(UAU^\dagger)^k=UA^kU^\dagger$.
Powers of a block-diagonal matrix are block diagonal, proving
\eqref{eq:exp-sum}. Powers of $I_a\otimes A\otimes I_c$ give
\eqref{eq:exp-id-tensor}. Each argument includes the zeroth power.
These are instances of the standard matrix-exponential calculus
\cite{higham-exp}.

For nonzero angles, Hermiticity of a direct sum implies Hermiticity of
every block, so \eqref{eq:exp-sum} preserves the typing premise.
In \eqref{eq:exp-id-tensor}, the right-hand rotation must be valid;
Hermiticity of the whole axis alone does not imply Hermiticity of $A$
when $a=0$ or $c=0$. Normalize that dimension-zero case to $I_0$ first.

There is no general law distributing an exponential over a tensor
product or a composition:
\[
 \Rot{A\otimes B}t\ne\Rot At\otimes\Rot Bt,
 \qquad
 \Rot{AB}t\ne\Rot At\Rot Bt.
\]
Both proposed laws already fail for $A=B=I_1$ and $t=2$: the left side
is $[-1]$ and the right side is $[1]$.
The familiar identity $e^{X+Y}=e^Xe^Y$ for commuting $X,Y$ involves
\emph{ordinary matrix addition}. Direct sum is different, and a sum of
two Hermitian unitaries need not be a valid UOp axis.

\begin{lstlisting}
splitRotationSum :: UOp -> Maybe UOp
splitRotationSum (R (DirectSum xs) t) =
  Just (directSum [rotation x t | x <- xs])
splitRotationSum _ = Nothing

conjugate :: UOp -> UOp -> UOp
conjugate u a = compose (dim a) [u, a, adjoint u]
\end{lstlisting}

\code{conjugate u a} assumes equal dimensions. Equation~\eqref{eq:exp-conjugate}
can move a conjugation into or out of a rotation depending on the
chosen optimization. Similarly, splitting a block rotation exposes
independent blocks; factoring equal-angle block rotations can make
the representation smaller. A rewrite strategy chooses the direction.

\section{Adjacent rewrites on composition lists}

Associativity makes every pair of adjacent factors directly accessible.
One generic list function is sufficient; there is no need to enumerate
the possible binary-tree shapes.

\begin{lstlisting}
type Rule = UOp -> Maybe UOp
type PairRule = UOp -> UOp -> Maybe UOp

rewriteAdjacent :: PairRule -> [UOp] -> Maybe [UOp]
rewriteAdjacent rule (a : b : rest) =
  case rule a b of
    Just c  -> Just (c : rest)
    Nothing -> (a :) <$> rewriteAdjacent rule (b : rest)
rewriteAdjacent _ _ = Nothing

inComposition :: PairRule -> Rule
inComposition rule (Compose d xs) =
  compose d <$> rewriteAdjacent rule xs
inComposition _ _ = Nothing
\end{lstlisting}

The first successful pair is replaced, preserving its prefix and
suffix. \code{Nothing} means that this rule did not match; it is not a
typing failure or a proof that no useful equality exists.

This deliberately simple scan makes one replacement per composition
per pass. Cancelling a list of $n$ copies of $X$ therefore takes
$\Theta(n)$ passes and $\Theta(n^2)$ total traversal work. A stack-based
scan can revisit newly adjacent factors within one traversal.

Same-axis rotation fusion, cancellation, and permutation composition
are small pair rules:

\begin{lstlisting}
fuse :: PairRule
fuse (R a s) (R b t)
  | a == b = Just (rotation a (s + t))
fuse (PermuteTensor left q) (PermuteTensor ds p)
  | left == reorder p ds =
      Just (permuteTensor ds (reorder q p))
fuse (PermuteSum left q) (PermuteSum ds p)
  | left == reorder p ds =
      Just (permuteSum ds (reorder q p))
fuse a b
  | a == adjoint b = Just (identity (dim a))
  | otherwise     = Nothing

fuseComposition :: Rule
fuseComposition = inComposition fuse
\end{lstlisting}

The permutation guard matches the \emph{output decomposition of the
right operand} with the input decomposition of the left operand.
Equal total dimensions alone do not justify composing their metadata.
All other pairs remain available to the enclosing composition.

Syntactic equality is sufficient but incomplete for these matches.
Adjoint comparison can also traverse an entire subtree; sharing or cached
adjoints can reduce that cost.

\subsection{Moving a permutation across a parallel layer}

Naturality also gives a useful adjacent-pair rewrite, without expanding
the permutation into basis indices:
\[
 P(\ds,p)\,\mathop{\star}_i A_i
 =\left(\mathop{\star}_j A_{p_j}\right)P(\ds,p),
 \qquad \dim A_i=d_i.
\]
Here $(P,\star)$ is either $(T,\otimes)$ or $(S,\oplus)$. This is the
conjugation law multiplied on the right by $P$. Notice that $P$ retains
its \emph{original input dimensions}; it is the parallel layer that is
reordered. The equality follows even when those dimensions are unequal.

\begin{lstlisting}
movePermutationRight :: PairRule
movePermutationRight p@(PermuteTensor ds order) (Tensor xs)
  | ds == map dim xs =
      Just (compose (product ds) [tensor (reorder order xs), p])
movePermutationRight p@(PermuteSum ds order) (DirectSum xs)
  | ds == map dim xs =
      Just (compose (sum ds) [directSum (reorder order xs), p])
movePermutationRight _ _ = Nothing

movePermutation :: Rule
movePermutation = inComposition movePermutationRight
\end{lstlisting}

The result contains a composition because the pair is replaced by
\emph{two} factors, not by one primitive. The enclosing composition
builder flattens it immediately. The same adjacent-pair scan used for
fusion therefore also handles this rule.

\paragraph{Matching shape is essential.}
The guard compares dimension \emph{lists}, not just their sums or
products. A permutation of blocks $[1,3]$ cannot be pushed through a
layer with blocks $[2,2]$ by reusing its index list. That would move
different basis coordinates. One must first establish a compatible
decomposition or leave the rule inapplicable.

Moving permutations right can expose fusion; moving them left can suit
a backend's output layout. Choose one direction for a pass.

\section{Interchange and composition lifting}

\subsection{Compose aligned factors or blocks}

When $\dim A_i=\dim B_i=d_i$ for every position,
\begin{align}
 \left(\bigotimes_i A_i\right)
 \left(\bigotimes_i B_i\right)
   &=\bigotimes_i(A_iB_i), \label{eq:tensor-interchange}\\
 \left(\bigoplus_i A_i\right)
 \left(\bigoplus_i B_i\right)
   &=\bigoplus_i(A_iB_i). \label{eq:sum-interchange}
\end{align}
The first is the mixed-product law for the Kronecker product
\cite{higham-kron}; the second follows by block multiplication.
Both preserve the left-to-right order inside each product $A_iB_i$.
There are no cross terms between distinct direct-sum blocks.

\begin{lstlisting}
interchange :: PairRule
interchange (Tensor as) (Tensor bs)
  | map dim as == map dim bs =
      Just (tensor (zipWith times as bs))
interchange (DirectSum as) (DirectSum bs)
  | map dim as == map dim bs =
      Just (directSum (zipWith times as bs))
interchange _ _ = Nothing

times :: UOp -> UOp -> UOp
times a b = compose (dim a) [a, b]

interchangeComposition :: Rule
interchangeComposition = inComposition interchange
\end{lstlisting}

Equality of the dimension lists checks both lengths and corresponding
dimensions, preventing \code{zipWith} from silently truncating a list.
The helper \code{times} assumes its operands have equal dimension.

For several aligned products, one could transpose a rectangular list
of operands and compose each column. Reusing the adjacent-pair rule is
usually simpler and already covers arbitrary composition lengths.
Different decompositions of the same total dimension do not match this
rule. For example, block sizes $[1,3]$ and $[2,2]$ cannot be zipped.
An arbitrary unitary block cannot be split at a desired boundary unless
it actually preserves that decomposition.

\paragraph{Useful consequences.}
Operators on disjoint tensor factors commute:
\[
 (A\otimes I_b)(I_a\otimes B)
   =A\otimes B
   =(I_a\otimes B)(A\otimes I_b),
 \qquad \dim A=a,\quad\dim B=b.
\]
This follows directly from \eqref{eq:tensor-interchange}; it does not
license arbitrary reordering of composition operands.
The two interchange laws also imply that tensor products and direct
sums of unitaries are unitary, by applying them to each operator and
its adjoint.

\subsection{Lift composition by padding and transposing}

\label{sec:lifting}

Orienting interchange outward gives a flat composition of parallel
layers, even when operand composition lists have different lengths. For example,
\begin{align*}
 (AB)\otimes C&=(A\otimes C)(B\otimes I_{\dim C}),\\
 (AB)\oplus C&=(A\oplus C)(B\oplus I_{\dim C}).
\end{align*}
The second layer contains an \emph{identity}, not another copy of $C$.
With two nontrivial lists the same construction gives
\[
 (A_1A_2)\otimes(B_1B_2B_3)
 =(A_1\otimes B_1)(A_2\otimes B_2)(I_{\dim A_1}\otimes B_3).
\]
These products retain mathematical order: the rightmost layer acts first.

\par\noindent\begin{minipage}{\linewidth}
\begin{proposition}[Composition lifting]
Let $\star$ denote either tensor product or direct sum. For each position
$j$, let $A_{j1},\ldots,A_{jm_j}$ be operators on $\C^{d_j}$, and let
$m=\max\{0,m_1,\ldots,m_k\}$. Extend each list to length $m$ by appending $I_{d_j}$.
Writing $\widehat A_{j\ell}$ for the padded entries, we have
\[
 \mathop{\star}_{j=1}^k\left(\prod_{\ell=1}^{m_j}A_{j\ell}\right)
 =\prod_{\ell=1}^m\left(\mathop{\star}_{j=1}^k\widehat A_{j\ell}\right).
\]
Every layer on the right has dimension $\prod_j d_j$ for tensor, or
$\sum_j d_j$ for direct sum. Empty products retain their annotated
dimension. When there are no operands or all products are empty, the
usual identity laws supply the result.
\end{proposition}
\begin{proof}
Padding does not change any operand. Apply the appropriate interchange
law repeatedly to the right-hand product of layers. In each position
the resulting product is the original operand followed by identities.
\end{proof}
\end{minipage}\par

\paragraph{The list layout.}
Regard each operand's factors as one row, pad the rows to a rectangle,
then transpose to obtain the layers:
\[
 \begin{bmatrix} A_1&A_2&I_a\\ B_1&B_2&B_3\end{bmatrix}
 \quad\longmapsto\quad
 \begin{bmatrix} A_1&B_1\\ A_2&B_2\\ I_a&B_3\end{bmatrix}.
\]
This produces $m$ layers with $k$ slots each, not a Cartesian product
of choices. Only identity slots are added; original factors are not
repeated. Repeated lifting through larger contexts may still duplicate
identity padding and increase the full tree's node count.

\subsection{One helper for both parallel operations}

The two cases share one helper. A non-composition is a one-element list;
an empty composition contributes identities of its own dimension.
The constructors passed to the helper are the structural builders above,
which do not themselves lift compositions.

\begin{lstlisting}
liftCompositions :: Rule
liftCompositions (Tensor xs) = liftParallel tensor product xs
liftCompositions (DirectSum xs) = liftParallel directSum sum xs
liftCompositions _ = Nothing

liftParallel :: ([UOp] -> UOp) -> ([Dim] -> Dim)
             -> [UOp] -> Maybe UOp
liftParallel con size xs
  | any isComposition xs =
      Just (compose (size (map dim xs)) (map con layers))
  | otherwise = Nothing
  where
    isComposition (Compose _ (_:_)) = True
    isComposition _                 = False
    factors (Compose _ ys) = ys
    factors a              = [a]
    rows = map factors xs
    height = maximum (0 : map length rows)
    pad a ys = ys ++ replicate (height - length ys)
                              (identity (dim a))
    layers = transpose (zipWith pad xs rows)
\end{lstlisting}

\paragraph{Padding preserves the decomposition.}
Use the original operand's dimension: an empty \code{Compose d []}
row still needs $I_d$. Pad before transposing so every layer keeps the
same factor or block dimensions. Transposing ragged lists loses slots;
cycling a short row repeats its operator.

\paragraph{Keep phases inside the chosen form.}
Scale one factor of a nonempty composition, as in
Section~\ref{sec:phases}. Wrapping it in
\code{Tensor [phase t,Compose d xs]} would reintroduce the nesting
that lifting removes.

\paragraph{Choose one orientation.}
\code{interchangeComposition} moves composition into a parallel
operation; \code{liftCompositions} moves it out. Use the former to
expose local cancellations and the latter to obtain layers. Installing
both in the same fixed point can make the expression oscillate.

\section{Tensor distribution}

\label{sec:distribution}

Leftmost direct sums distribute literally in the chosen basis order:
\begin{equation}
 \left(\bigoplus_i B_i\right)\otimes C
   =\bigoplus_i(B_i\otimes C). \label{eq:left-distribute}
\end{equation}
This follows immediately from the block definition of the Kronecker
product. The corresponding list rule is a map over the branches:

\begin{lstlisting}
distributeFirst :: Rule
distributeFirst (Tensor (DirectSum branches : rest)) =
  Just (directSum [tensor (b : rest) | b <- branches])
distributeFirst _ = Nothing
\end{lstlisting}

Distribution over another operand must account for the basis order.
For instance,
\[
 I_2\otimes([1]\oplus[-1])=\diag(1,-1,1,-1),\qquad
 (I_2\otimes[1])\oplus(I_2\otimes[-1])=\diag(1,1,-1,-1).
\]
They are permutation-conjugate, not equal. A rewrite that merely maps
the surrounding tensor factors into each branch changes the operator.

To distribute a selected sum, write
\[
 M=A\otimes\left(\bigoplus_i B_i\right)\otimes C,
 \qquad a=\dim A,\quad b=\sum_i\dim B_i,\quad c=\dim C.
\]
Set $P=T([a,b,c],[1,0,2])$, which moves the sum before $A$.
Tensor naturality and \eqref{eq:left-distribute} give
\[
 PMP^\dagger
  =\left(\bigoplus_i B_i\right)\otimes A\otimes C
  =\bigoplus_i(B_i\otimes A\otimes C).
\]
Thus the exact matrix-preserving rewrite is
\begin{equation}
 \boxed{\displaystyle
 M=P^\dagger\left(\bigoplus_i(B_i\otimes A\otimes C)\right)P.}
 \label{eq:distribute}
\end{equation}
It is especially simple because it allows $B_i$ to appear first inside
each output block. No basis permutation of dimension $abc$ is expanded;
$P$ stores just three dimensions and three indices.

If the desired block order is instead $A\otimes B_i\otimes C$, put
$b_i=\dim B_i$ and define
\[
 D=\left[\left(\bigoplus_i T([b_i,a],[1,0])\right)
               T([a,b],[1,0])\right]\otimes I_c.
\]
Then $DMD^\dagger=\bigoplus_i(A\otimes B_i\otimes C)$, so
$M=D^\dagger(\bigoplus_i(A\otimes B_i\otimes C))D$.
This version preserves the original factor order inside each branch,
at the cost of additional symbolic shuffles.

\Needspace{30\baselineskip}
\subsection{Match a prefix, a selected sum, and a suffix}

The general rule needs only one list split. \code{break} finds the
first direct-sum operand; \code{before} and \code{after} supply $A$
and $C$ in \eqref{eq:distribute}.

\begin{lstlisting}
distributeTensor :: Rule
distributeTensor (Tensor operands) =
  case break isDirectSum operands of
    (before, DirectSum branches : after) ->
      let a = product (map dim before)
          b = sum (map dim branches)
          c = product (map dim after)
          shuffle = permuteTensor [a, b, c] [1, 0, 2]
          expanded = directSum
            [ tensor (branch : (before ++ after))
            | branch <- branches
            ]
      in Just (compose (a * b * c)
           [adjoint shuffle, expanded, shuffle])
    _ -> Nothing
distributeTensor _ = Nothing

isDirectSum :: UOp -> Bool
isDirectSum (DirectSum _) = True
isDirectSum _             = False
\end{lstlisting}

Empty prefixes or suffixes have dimension one. An empty branch list
has dimension zero. The formula remains valid in each case, although
normalization usually eliminates empty sums and other zero-dimensional
terms before this rule runs.

Try \code{distributeFirst} first: it avoids an identity shuffle that
the minimal permutation constructor may not recognize when a factor
has dimension one.

\paragraph{Factoring is the inverse rule, with the same ordering premises.}
For example, a common tensor suffix can be factored exactly:
\[
 \bigoplus_i(B_i\otimes C)
   =\left(\bigoplus_i B_i\right)\otimes C.
\]
Matching explicit tensor suffixes requires only comparing lists, not
searching all binary parenthesizations. A common tensor prefix instead
needs the corresponding shuffle. Do not discard it when reversing
\eqref{eq:distribute}.

Distribution can duplicate the surrounding factors into many branches.
Keep it out of the structural builders and select expansions using the
budgets in Section~\ref{sec:budgets}.

\Needspace{32\baselineskip}
\section{Traversal, strategy, and soundness}

\subsection{Traversal and cleanup}

\label{sec:traversal}

A rule inspects one root. A traversal decides where roots are visited.
The following two helpers work for every local rule in this note:

\begin{lstlisting}
mapChildren :: (UOp -> UOp) -> UOp -> UOp
mapChildren f op = case op of
  R a t        -> R (f a) t
  Tensor xs    -> Tensor (map f xs)
  DirectSum xs -> DirectSum (map f xs)
  Compose d xs -> Compose d (map f xs)
  _            -> op

bottomUp :: (UOp -> UOp) -> UOp -> UOp
bottomUp f = go
  where go a = f (mapChildren go a)

cleanup :: UOp -> UOp
cleanup op = case op of
  R a t             -> rotation a t
  Tensor xs         -> tensor xs
  DirectSum xs      -> directSum xs
  Compose d xs      -> compose d xs
  PermuteTensor ds p -> permuteTensor ds p
  PermuteSum ds p    -> permuteSum ds p
\end{lstlisting}

\code{mapChildren} preserves the root and metadata; \code{cleanup}
rebuilds through the structural constructors. Bottom-up traversal
processes children first, including rotation axes. A replacement can
introduce new redexes, so one traversal need not reach a fixed point.

\subsection{Rule priority and stopping}

\label{sec:strategy}

Rule priority is explicit. Each local call chooses the first applicable
rule; a later traversal can revisit the result. Here \code{<|>} chooses
the first \code{Just}, and \code{fromMaybe} leaves an unmatched term alone.

\begin{lstlisting}
firstRule :: [Rule] -> Rule
firstRule [] _ = Nothing
firstRule (r:rs) a = r a <|> firstRule rs a

applyRules :: [Rule] -> UOp -> UOp
applyRules rules a = fromMaybe a (firstRule rules a)

pass :: [Rule] -> UOp -> UOp
pass rules = bottomUp (cleanup . applyRules rules . cleanup)

flatPass :: UOp -> UOp
flatPass = pass [fuseComposition, liftCompositions]

data RewriteResult = Stable UOp | FuelExhausted UOp
  deriving (Show, Eq)

repeatPass :: Int -> (UOp -> UOp) -> UOp -> RewriteResult
repeatPass fuel step a
  | fuel <= 0 = FuelExhausted a
  | b == a    = Stable a
  | otherwise = repeatPass (fuel - 1) step b
  where b = step a
\end{lstlisting}

\code{flatPass} chooses outward interchange. Cleanup before matching
and after replacement removes incidental units and associations.

\paragraph{An equality test is not a termination proof.}
A fuel bound makes the driver terminate even if a proposed rule set
cycles. \code{FuelExhausted} returns a sound intermediate term, not a
claimed normal form. \code{Stable} means that this particular composite
pass made no syntactic change; it does not establish semantic canonicity
or exclude redexes for rules that the strategy did not select. For an
unbounded driver, justify termination of the \emph{whole} oriented rule
set using a well-founded measure. Soundness of its individual equations
does not provide such a measure.

\subsection{Expansion budgets}

\label{sec:budgets}

Distribution across $r$ sums with $k_1,\ldots,k_r$ branches can expose
$\prod_j k_j$ combinations. Its growth needs a separate strategy decision
from composition lifting, even though both transformations are symbolic.

A small wrapper bounds the number of nodes in each replacement subtree.
It does not bound the size of the enclosing expression:

\begin{lstlisting}
nodeCount :: UOp -> Integer
nodeCount op = 1 + case op of
  R a _        -> nodeCount a
  Tensor xs    -> sum (map nodeCount xs)
  DirectSum xs -> sum (map nodeCount xs)
  Compose _ xs -> sum (map nodeCount xs)
  _            -> 0

limitReplacementNodes :: Integer -> Rule -> Rule
limitReplacementNodes limit rule a = case rule a of
  Just b | nodeCount b <= limit -> Just b
  _ -> Nothing
\end{lstlisting}

Applying this local guard at several positions can still exceed a
whole-expression budget. For a tree with $N$ nodes, replacing an
occurrence of $a$ by $b$ changes the total to
$N-\code{nodeCount}\ a+\code{nodeCount}\ b$; a global budget belongs
in the strategy that knows that context.

The guard also constructs and visits the candidate before accepting it.
Before a large expansion, estimate its branch count and duplicated
context size from the matched lists, and reject it before allocation
if necessary. A separate pass or rewrite bound controls search effort.

Choose a cost appropriate to the task: tree size, shared nodes, depth,
branch count, or backend operations. Flat layers can improve scheduling
while adding identities; a factored term can preserve a smaller
description. Run expansion and factoring as separate directed phases,
and report a budget-limited result as partial.

\subsection{Contextual soundness}

\par\noindent\begin{minipage}{\linewidth}
\begin{proposition}[Soundness of rewriting]
Suppose a well-formed term $a$ is replaced by $b$ using any algebraic
rule above, with its stated premises satisfied.
Then $b$ is a valid term of the same dimension and
$\sem b=\sem a$. The replacement remains valid inside any well-typed
expression context, including a rotation axis.
\end{proposition}
\begin{proof}
Each rule is one of the proved matrix identities, with matching
dimensions. For newly introduced rotations, Hermiticity follows
from the stated premise, restriction to a direct-sum block, or unitary
conjugation. Existing unitarity is preserved by equality. Every
constructor's denotation is a function of its children's denotations,
so replacement by an equal matrix preserves the enclosing denotation.
In a rotation axis it also preserves Hermiticity. Induction gives
the result for a sequence of rewrites.
\end{proof}
\end{minipage}\par

This is a sound collection of rules, not a complete procedure for
semantic equality or a claim of a unique algebraic normal form.
Lists remove associativity from pattern matching; they do not eliminate
the choices between factored and expanded forms, block decompositions,
or equivalent rotation axes.

\section{QFT example and implementation checks}

The QOp generator \code{qft n} acts on $n$ qubits, hence dimension $2^n$.
Lowering and normalization operate on its syntax. In the repository:

\begin{lstlisting}[caption={Repository entry point; separate from the standalone listings.}]
import Programs.QFT (qft)
import qualified QPP.UOp as U

qftUOp n = fmap U.normalize (U.fromQOp (qft n))
\end{lstlisting}

The Hadamard lowering is
\[
 H=\Ph{1/2}\otimes
   \bigl(R_Z(1/2)R_X(1/2)R_Z(1/2)\bigr),\quad
 X=S([1,1],[1,0]),\quad Z=I_1\oplus\Ph1.
\]
After splitting the diagonal rotations and placing the scalar phase
on the first factor, write
\[
 D_+=\Ph{1/4}\oplus\Ph{3/4},\qquad
 D_-=\Ph{-1/4}\oplus\Ph{1/4},\qquad K=R_X(1/2).
\]
Then $H=D_+KD_-$ exactly, including its global phase. In a larger register,
\[
 I_a\otimes H\otimes I_c
 =(I_a\otimes D_+\otimes I_c)
  (I_a\otimes K\otimes I_c)
  (I_a\otimes D_-\otimes I_c).
\]
This is the lifting rule applied to the actual three-factor lowering.
For a controlled relative phase, the phase laws similarly give
$I_2\oplus(\Ph q\otimes R_Z(2q))=I_3\oplus\Ph{2q}$.
With intervening qubits, the appropriate tensor identity padding remains.

\begin{center}
\begin{tabular}{@{}lrrrrr@{}}
\toprule
& Lowered nodes & Lowered depth & Flat nodes & Flat depth & Factors\\
\midrule
\code{qft 4}  & 154 & 12 & 127  & 7 & 19\\
\code{qft 16} & 1876 & 36 & 1561 & 7 & 169\\
\bottomrule
\end{tabular}
\end{center}
These are constructor-occurrence counts for the current implementation,
not gate counts or claims of minimum size. Each Hadamard contributes
three factors, each of the $n(n-1)/2$ controlled phases one, and the
final register reversal one. For $n\ge2$ this gives
$3n+n(n-1)/2+1$ outer factors.

For \code{qft 4}, structural simplification before lifting used 110 nodes
at depth 10; lifting gives 127 at depth 7. The additional nodes are
identity padding and layer structure. The payoff is one outer
\code{Compose 16} with 19 factors, each ready for local list-based rules.
The files under \code{output/qft4/} and \code{output/qft16/} record both
raw and normalized terms; normalized pretty prints use one factor per line.

\Needspace{38\baselineskip}
\subsection{Repository implementation and checks}

\label{sec:implementation}

\code{QPP.UOp} exposes syntax, pretty printing, and builders that flatten
associations. Its production \code{normalize} also collects phases,
simplifies rotations, cancels or fuses adjacent factors, and lifts
compositions out of tensor products and direct sums. It uses a stack-based
adjacent scan and repeats its specialized pass until syntactic equality.
The result has no nonempty \code{Compose} immediately inside a tensor
or sum; rotations may still have composite axes.

Tensor products of sums remain factored. The distribution, permutation
movement, and generic fuel/budget strategies in this note are proposed
patterns. General UOp validation is also future work: the normalizer
expects valid input, for example from \code{fromQOp}. Explicit proof
objects are optional. Dimensions or adjoints can be cached if repeated
tree walks become costly.

Lowering and normalization operate symbolically. Their tests use an
independent interpreter that applies UOp terms directly to small
state vectors. A rotation acts by
\[
 R_A(t)v=\cos(\pi t/2)v-i\sin(\pi t/2)Av,
\]
and tensors, sums, compositions, and permutations act directly on vector
slices or coordinates. Complex, nonuniform amplitudes expose phase and
basis-order errors that an all-ones vector can miss.

Separate small-instance checks compare raw and normalized lowering
against QOp's matrix semantics on every basis vector. These reference
matrices are confined to tests, including QFT instances up to three
qubits; the lowering and normalization algorithms construct none.
Both kinds of finite numerical checks are regression evidence, while
the equations supply the general correctness arguments.

Tests for these rules should cover:
\begin{enumerate}
\item \textbf{Permutations:} non-involutive orders, unequal and zero-sized
blocks, adjoints, fusion, and the non-Hermitian $[1,2]$ block swap.
\item \textbf{Distribution:} nonempty prefixes and suffixes, unequal or
zero-sized branches, and non-diagonal block operators.
\item \textbf{Interchange and lifting:} mismatched decompositions,
unequal or empty composition lists, dimension-correct padding, and
an outer composition without repeated operands or phases.
\item \textbf{Rotations:} adjoints, fusion, block splitting, conjugation,
and the original-angle distinction between $0$ and $4$ for a
non-Hermitian axis.
\item \textbf{Invariants:} dimension, denotation, idempotence, and
Hermiticity inside rewritten axes; proof-object preservation if used.
\end{enumerate}

The uncaptioned listings form one Haskell module using only \code{base},
with the imports recorded at the top of the \LaTeX{} source. The captioned
QFT listing imports the repository. The reference implementation relies
on the validity contract in Section~\ref{sec:contracts}.

The production suite checks vector action and idempotence on 500
generated terms, plus QFT shape and dimension through \code{qft 16}.
QFT vector checks stop at four qubits; reference-matrix checks stop at
three. New expansion rules need their own regression tests when added.

\begin{thebibliography}{9}
\bibitem{higham-kron}
Nicholas J. Higham.
\emph{What Is the Kronecker Product?} 25 August 2020.
\href{https://nhigham.com/2020/08/25/what-is-the-kronecker-product/}
{nhigham.com: Kronecker product}.
Background for the mixed-product and adjoint identities.

\bibitem{higham-exp}
Nicholas J. Higham.
\emph{What Is the Matrix Exponential?} 28 May 2020.
\href{https://nhigham.com/2020/05/28/what-is-the-matrix-exponential/}
{nhigham.com: matrix exponential}.
Background for the power-series and commuting-exponential identities.
\end{thebibliography}
\end{document}